\vskip -10pt
In this paper we draw a connection between causal inference with proxy variables and the ground-breaking work in the machine learning community on latent variable models. Since almost all observational studies rely on proxy variables, this connection is highly relevant.

We introduce a model which is the first attempt at tying these two ideas together: The Causal Effect Variational Autoencoder (CEVAE), a neural network latent variable model used for estimating individual and population causal effects. In extensive experiments we showed that it is competitive with the state-of-the art on benchmark datasets, and more robust to hidden confounding both at a toy artificial dataset as well as modifications of real datasets, such as the newly introduced Twins dataset. 
For future work, we plan to employ the expanding set of tools available for latent variables models (e.g. \citet{kingma2016improving,tran2015variational,maaloe2016auxiliary,ranganath2016operator}), as well as to further explore connections between method of moments approaches such as \citet{anandkumar2014tensor} with the methods for effect restoration given by \citet{kuroki2014measurement,miao2016identifying}. 

